\section{Разложение многочленов на множители}

Следующие результаты помогают решить, приводим ли многочлен, и найти его множители.

\subsection*{Рациональные корни}

\begin{proposition}[1.11]
Пусть $r\in\Q$~--- корень многочлена
\[
f=a_mX^m+a_{m-1}X^{m-1}+\dots+a_0,\qquad a_i\in\Z,
\]
и $r=c/d$, где $c,d\in\Z$, $\gcd(c,d)=1$. Тогда $c\mid a_0$ и $d\mid a_m$.
\end{proposition}

\begin{proof}
Умножив равенство $f(c/d)=0$ на $d^m$, получаем
\[
a_mc^m+a_{m-1}c^{m-1}d+\dots+a_1cd^{m-1}+a_0d^m=0 .
\]
Все слагаемые, кроме первого, делятся на $d$, поэтому $d\mid a_mc^m$. Так как $\gcd(d,c^m)=1$ (у $d$ и $c$ нет общих
простых делителей), то $d\mid a_m$. Аналогично все слагаемые, кроме последнего, делятся на $c$, поэтому $c\mid a_0d^m$
и $c\mid a_0$.
\end{proof}

В частности, \emph{рациональный корень унитарного многочлена с целыми коэффициентами~--- целое число} ($d\mid1$).

\begin{example}[1.12]
Многочлен $f(X)=X^3-3X-1$ неприводим в $\Q[X]$. Если бы многочлен степени $3$ раскладывался в произведение многочленов
меньшей степени, один из множителей был бы линейным, и $f$ имел бы рациональный корень. По \rf{1.11} возможные
рациональные корни~--- только $\pm1$, а $f(1)=-3\ne0$, $f(-1)=1\ne0$.
\end{example}

То же рассуждение годится для любого многочлена степени $2$ или $3$ над любым полем: он приводим тогда и только тогда,
когда имеет корень. Для степени $4$ это уже неверно: многочлен может быть произведением двух неприводимых квадратных
(см.\ ниже \rf{1.18d}).

\subsection*{Лемма Гаусса}

Для $p$ простого через $f\mapsto\bar f$ обозначим \emph{приведение по модулю $p$}: гомоморфизм колец $\Z[X]\to\F_p[X]$,
заменяющий каждый коэффициент его вычетом по модулю $p$. Он сюръективен, а его ядро~--- многочлены, все коэффициенты
которых делятся на $p$, т.\,е.\ $p\Z[X]$; по теореме о гомоморфизме $\Z[X]/p\Z[X]\cong\F_p[X]$.

\par\noindent\begin{minipage}{\linewidth}
\[
\begin{tikzcd}[column sep=large,row sep=large]
\Z[X]\arrow[r,"f\,\mapsto\,\bar f",cold,two heads]\arrow[d,two heads] & \F_p[X]\\
\Z[X]/p\Z[X]\arrow[ur,cnew,"\sim"'] &
\end{tikzcd}
\]
\figcap{Приведение по модулю $p$ сюръективно и имеет ядро $p\Z[X]$, поэтому индуцирует изоморфизм
$\Z[X]/p\Z[X]\cong\F_p[X]$ (красная стрелка).}
\end{minipage}\par\medskip

\begin{proposition}[1.13: лемма Гаусса]
Пусть $f(X)\in\Z[X]$. Если $f(X)$ раскладывается нетривиально в $\Q[X]$ (в произведение многочленов меньшей степени),
то он раскладывается нетривиально и в $\Z[X]$.
\end{proposition}

\begin{proof}
Пусть $f=gh$ в $\Q[X]$, где $g,h$ непостоянны. Для подходящих целых $m,n\ge1$ (общих знаменателей) многочлены
$g_1=mg$ и $h_1=nh$ имеют целые коэффициенты, и мы получаем разложение
\[
mnf=g_1\cdot h_1\quad\text{в }\Z[X].
\]
Пусть простое $p$ делит $mn$. Приводя по модулю $p$, получаем
\[
0=\bar g_1\cdot\bar h_1\quad\text{в }\F_p[X].
\]
Так как $\F_p[X]$~--- область целостности, $\bar g_1=0$ или $\bar h_1=0$, т.\,е.\ $p$ делит все коэффициенты одного из
многочленов $g_1,h_1$, скажем $g_1$: $g_1=pg_2$, $g_2\in\Z[X]$. Получаем разложение
\[
(mn/p)f=g_2\cdot h_1\quad\text{в }\Z[X].
\]
Продолжая так, мы удалим все простые множители $mn$ и получим разложение $f=G\cdot H$ в $\Z[X]$, где $G=ug$, $H=vh$,
$u,v\in\Q^\times$. Степени при умножении на ненулевые константы не меняются, так что разложение нетривиально.
\end{proof}

Доказательство даёт больше: \emph{множители $G,H\in\Z[X]$ отличаются от исходных $g,h$ лишь постоянными множителями},
$G=ug$, $H=u^{-1}h$ ($uv=1$, так как $GH=f=gh$).

\begin{proposition}[1.14]
Если $f\in\Z[X]$ унитарен, то каждый унитарный делитель $f$ в $\Q[X]$ лежит в $\Z[X]$.
\end{proposition}

\begin{proof}
Пусть $g$~--- унитарный делитель $f$ в $\Q[X]$: $f=gh$, где $h\in\Q[X]$ тоже унитарен. Выберем натуральные $m,n$ с
наименьшим числом простых множителей, для которых $mg,nh\in\Z[X]$. Как в доказательстве леммы Гаусса, если простое $p$
делит $mn$, то оно делит все коэффициенты одного из многочленов $mg$, $nh$, скажем $mg$. В частности, $p$ делит его
старший коэффициент $m$ ($g$ унитарен). Но тогда $\frac mp g\in\Z[X]$, что противоречит выбору $m$. Значит, $mn=1$,
$m=n=1$ и $g,h\in\Z[X]$.
\end{proof}

\begin{remark}[1.14a]
Утверждение \rf{1.14}~--- это лемма Гаусса для унитарных многочленов, и его можно вывести из \rf{1.13}. По замечанию после
доказательства \rf{1.13} имеем $f=(ug)(u^{-1}h)$ с $ug,u^{-1}h\in\Z[X]$. Их старшие коэффициенты $u$ и $u^{-1}$~--- целые
числа, поэтому $u=\pm1$ и $g=\pm ug\in\Z[X]$. Без унитарности $g$ утверждение неверно:
$X^2-1=(2X-2)\cdot\tfrac12(X+1)$. Само по себе \rf{1.13} не утверждает, что \emph{данные} множители $g,h$ целые,~--- только
что после умножения на подходящие константы они становятся целыми.

Если же заранее известно, что унитарный делитель $g$ лежит в $\Z[X]$, то лемма Гаусса не нужна: \emph{если $f,g\in\Z[X]$
унитарны и $f=gh$ с $h\in\CC[X]$, то $h\in\Z[X]$}. Действительно, на унитарный многочлен можно делить с остатком в $\Z[X]$
(в доказательстве \rf{1.6} делить приходится только на старший коэффициент $g$, равный $1$): $f=gq+r$, $q,r\in\Z[X]$,
$\deg r<\deg g$. По единственности деления с остатком в $\CC[X]$ получаем $h=q$, $r=0$.
\end{remark}

\begin{aside}[1.15]
Вот другой взгляд на \rf{1.14}. Комплексное число $\alpha$ называется \emph{целым алгебраическим}, если оно~--- корень
унитарного многочлена из $\Z[X]$. По \rf{1.11} все целые алгебраические числа, лежащие в $\Q$, лежат в $\Z$. Целые
алгебраические числа образуют подкольцо в $\CC$ (это доказывается в коммутативной алгебре). Пусть теперь
$\alpha_1,\dots,\alpha_m$~--- корни $f$ в $\CC$; по определению это целые алгебраические числа. Коэффициенты любого
унитарного делителя $f$~--- многочлены с целыми коэффициентами от некоторых из $\alpha_i$ (с точностью до знака это
элементарные симметрические функции от его корней), и потому целые алгебраические. Если они рациональны, то они целые.

В частности, \emph{минимальный многочлен над $\Q$ целого алгебраического числа $\alpha$ лежит в $\Z[X]$}: он унитарен и
делит в $\Q[X]$ унитарный многочлен из $\Z[X]$ с корнем $\alpha$ (минимальный многочлен и это его свойство~--- ниже,
\rf{1.27c}(i)), так что можно применить \rf{1.14}.
\end{aside}

\subsection*{Критерий Эйзенштейна}

\begin{proposition}[1.16: критерий Эйзенштейна]
Пусть
\[
f=a_mX^m+a_{m-1}X^{m-1}+\dots+a_0,\qquad a_i\in\Z,
\]
и существует простое число $p$, для которого:
\begin{itemize}
\item $p$ не делит $a_m$;
\item $p$ делит $a_{m-1},\dots,a_0$;
\item $p^2$ не делит $a_0$.
\end{itemize}
Тогда $f$ неприводим в $\Q[X]$.
\end{proposition}

\begin{proof}
Если $f$ раскладывается нетривиально в $\Q[X]$, то по \rf{1.13} и в $\Z[X]$:
\[
a_mX^m+a_{m-1}X^{m-1}+\dots+a_0=(b_rX^r+\dots+b_0)(c_sX^s+\dots+c_0),
\]
где $b_i,c_i\in\Z$, $b_rc_s\ne0$, $r,s\ge1$, $r+s=m$. Так как $p$ делит $a_0=b_0c_0$, а $p^2$ не делит, $p$ делит ровно одно из
чисел $b_0$, $c_0$; пусть $p\mid b_0$, $p\nmid c_0$. Докажем индукцией, что $p\mid b_i$ при $0\le i\le r$. Пусть
$1\le i\le r$ и $p\mid b_0,\dots,b_{i-1}$. Так как $i\le r<m$, имеем $p\mid a_i$, и из равенства
\[
a_i=b_0c_i+b_1c_{i-1}+\dots+b_{i-1}c_1+b_ic_0
\]
(где $c_j=0$ при $j>s$) получаем $p\mid b_ic_0$, а значит, $p\mid b_i$. В частности, $p\mid b_r$, и $p$ делит
$a_m=b_rc_s$~--- противоречие.
\end{proof}

\begin{remark}[1.16a]
Критерий говорит о неприводимости в $\Q[X]$; в $\Z[X]$ многочлен может разложиться за счёт постоянного множителя:
$2X+6=2(X+3)$ удовлетворяет условиям при $p=3$, но в $\Z[X]$ это нетривиальное разложение ($2$ необратим в $\Z$).

Критерий часто применяют после замены переменной: многочлен $f(X)$ неприводим тогда и только тогда, когда неприводим
$f(X+c)$, так как $X\mapsto X+c$~--- автоморфизм кольца $F[X]$, сохраняющий степени.
\begin{enumerate}[label=(\alph*)]
\item $X^n-p$ ($p$ простое, $n\ge1$) неприводим над $\Q$. Поэтому над $\Q$ есть неприводимые многочлены любой степени.
\item $X^4+1$ неприводим над $\Q$: $(X+1)^4+1=X^4+4X^3+6X^2+4X+2$~--- многочлен Эйзенштейна при $p=2$.
\item Для простого $p$ многочлен $\Phi_p(X)=\dfrac{X^p-1}{X-1}=X^{p-1}+\dots+X+1$ неприводим над $\Q$:
\[
\Phi_p(X+1)=\frac{(X+1)^p-1}X=X^{p-1}+\binom p{p-1}X^{p-2}+\dots+\binom p2X+\binom p1,
\]
и все коэффициенты, кроме старшего, делятся на $p$ (см.\ раздел о Фробениусе), а свободный член $p$ не делится на $p^2$. (Это используется в \rf{1.42}.)
\end{enumerate}
\end{remark}

\subsection*{Многочлены над факториальным кольцом}

Утверждения \rf{1.13}, \rf{1.14} и \rf{1.16} верны с очевидными изменениями, если заменить $\Z$ факториальным
кольцом $R$, $\Q$~--- полем частных $K$ кольца $R$, а простое число $p$~--- простым (= неприводимым, \rf{1.0b})
элементом $\pi$ кольца $R$. Приведём
аккуратные формулировки и доказательства.

Наибольший общий делитель в $R$ определён с точностью до обратимого множителя (его можно найти по разложениям на простые
множители), и $\gcd(ab_1,\dots,ab_k)=a\gcd(b_1,\dots,b_k)$. \emph{Содержание} $c(f)$ ненулевого многочлена $f\in R[X]$~---
наибольший общий делитель его коэффициентов; $f$ \emph{примитивен}, если $c(f)$ обратимо, т.\,е.\ ни один простой элемент
не делит все коэффициенты $f$. Каждый ненулевой $f\in K[X]$ записывается как $f=cf_0$ с $c\in K^\times$ и примитивным
$f_0\in R[X]$: если $0\ne e\in R$~--- общий знаменатель коэффициентов $f$, то $ef\in R[X]$, и $ef=c(ef)\,f_0$. Такая запись
единственна с точностью до обратимых элементов $R$: если $cf_0=c'f_0'$, то, умножив на общий знаменатель, получим
$af_0=a'f_0'$ с $a,a'\in R$, и сравнение содержаний даёт $a=ua'$ с обратимым $u$.

\begin{lemma}[1.16b: лемма Гаусса для факториальных колец]
Произведение примитивных многочленов из $R[X]$ примитивно. Следовательно, $c(gh)=c(g)c(h)$ (с точностью до обратимого
множителя) для ненулевых $g,h\in R[X]$.
\end{lemma}

\begin{proof}
Пусть $g,h$ примитивны, а простой элемент $\pi$ делит все коэффициенты $gh$. Кольцо $R/(\pi)$~--- область целостности
(идеал $(\pi)$ прост), поэтому $(R/(\pi))[X]$~--- тоже область целостности (старший коэффициент произведения равен
произведению старших коэффициентов). Приводя по модулю $\pi$, получаем $\bar g\bar h=0$, откуда $\bar g=0$ или $\bar h=0$,
т.\,е.\ $\pi$ делит все коэффициенты $g$ или $h$~--- противоречие. Второе утверждение: $g=c(g)g_0$, $h=c(h)h_0$ с
примитивными $g_0,h_0$, и $gh=c(g)c(h)\,g_0h_0$, где $g_0h_0$ примитивен.
\end{proof}

\begin{theorem}[1.16c]
Пусть $R$~--- факториальное кольцо, $K$~--- его поле частных, $f\in R[X]$.
\begin{enumerate}[label=(\alph*)]
\item Если $f=gh$ в $K[X]$, то $f=GH$, где $G=ug\in R[X]$, $H=u^{-1}h\in R[X]$, $u\in K^\times$. В частности, если $f$
раскладывается в $K[X]$ на множители меньшей степени, то и в $R[X]$.
\item Если $f$ унитарен, то каждый его унитарный делитель в $K[X]$ лежит в $R[X]$.
\item \emph{(Критерий Эйзенштейна.)} Пусть $f=a_mX^m+\dots+a_0$, $m\ge1$, и $\pi$~--- простой элемент $R$, для которого
$\pi\nmid a_m$, $\pi\mid a_i$ при $i<m$ и $\pi^2\nmid a_0$. Тогда $f$ неприводим в $K[X]$.
\end{enumerate}
\end{theorem}

\begin{proof}
(a) Запишем $g=cg_0$, $h=c'h_0$ с $c,c'\in K^\times$ и примитивными $g_0,h_0$. По \rf{1.16b} многочлен $g_0h_0$ примитивен,
и $f=(cc')g_0h_0$. Запишем $cc'=a/b$ с $a,b\in R$ без общих простых делителей. Тогда $bf=a\,g_0h_0$; содержание левой
части делится на $b$, а содержание правой равно $a$ (с точностью до обратимого). Значит, $b\mid a$, т.\,е.\ $b$ обратим и
$cc'\in R$. Положим $G=cc'g_0=c'g$ и $H=h_0=c'^{-1}h$.

(b) Пусть $f=gh$, $g$ унитарен; тогда $h$ тоже унитарен. По (a) $f=(ug)(u^{-1}h)$ с $ug,u^{-1}h\in R[X]$; их старшие
коэффициенты $u$ и $u^{-1}$ лежат в $R$, т.\,е.\ $u$ обратим в $R$, и $g=u^{-1}(ug)\in R[X]$.

(c) Дословно повторяет доказательство \rf{1.16}: по (a) возникает разложение $f=(b_rX^r+\dots+b_0)(c_sX^s+\dots+c_0)$ в
$R[X]$ с $r,s\ge1$; так как $\pi$ прост, а $\pi^2\nmid a_0=b_0c_0$, он делит ровно одно из $b_0,c_0$, скажем $b_0$; далее
из $\pi\mid b_ic_0$, $\pi\nmid c_0$ следует $\pi\mid b_i$ (простота $\pi$), и в итоге $\pi\mid b_rc_s=a_m$~---
противоречие.
\end{proof}

Из (a) и \rf{1.16b} выводится также, что кольцо $R[X]$ факториально (и потому факториальны $\Z[X_1,\dots,X_n]$ и
$F[X_1,\dots,X_n]$ для поля $F$); нам это не понадобится.

\begin{example}[1.16d]
Пусть $k$~--- поле, $t$~--- переменная. Кольцо $R=k[t]$ факториально, $K=k(t)$, а $t$~--- простой элемент $R$.
\begin{enumerate}[label=(\alph*)]
\item $X^n-t$ неприводим над $k(t)$ при любом $n\ge1$.
\item При $\operatorname{char}k=p$ многочлен $X^{2p}+tX^p+t$ неприводим над $k(t)$ (коэффициенты при младших степенях
$t$, $t$ и $0$ делятся на $t$, свободный член не делится на $t^2$). При этом его производная равна $0$, и у него всего
два различных корня, каждый кратности $p$ (см.\ пример \rf{2.20a}).
\end{enumerate}
\end{example}

\subsection*{Алгоритм разложения в $\Q[X]$}

\begin{remark}[1.17]
Существует алгоритм разложения многочлена из $\Q[X]$ на неприводимые множители. Покажем это.

\emph{Шаг 1: сведение к унитарным целочисленным многочленам.} Пусть $f\in\Q[X]$, $\deg f=m$. Умножив $f$ на рациональное
число, сделаем его унитарным: $f=X^m+c_1X^{m-1}+\dots+c_m$. Пусть $D$~--- общий знаменатель $c_i$. Тогда
\[
D^mf\Bigl(\frac XD\Bigr)=X^m+Dc_1X^{m-1}+D^2c_2X^{m-2}+\dots+D^mc_m
\]
унитарен и лежит в $\Z[X]$, а разложения $f$ и этого многочлена соответствуют друг другу: если $f=gh$, $\deg g=k$, то
$D^mf(X/D)=D^kg(X/D)\cdot D^{m-k}h(X/D)$, и обратно. Например, $X^2+\frac12X+\frac13$ при $D=6$ превращается в
$X^2+3X+12$. Итак, достаточно рассматривать многочлены
\[
f(X)=X^m+a_1X^{m-1}+\dots+a_m,\qquad a_i\in\Z .
\]

\emph{Шаг 2: корни ограничены.} По основной теореме алгебры (\rf{5.6}) $f$ раскладывается в $\CC[X]$ на линейные множители:
$f(X)=\prod_{i=1}^m(X-\alpha_i)$, $\alpha_i\in\CC$. Из равенства
\[
0=f(\alpha_i)=\alpha_i^m+a_1\alpha_i^{m-1}+\dots+a_m
\]
следует, что $|\alpha_i|$ не превосходит числа, зависящего только от степени и коэффициентов $f$ (\rf{1.17a} ниже):
\[
|\alpha_i|\le\max\{1,mB\},\qquad B=\max_i|a_i| .
\]

\emph{Шаг 3: перебор.} Если $g$~--- унитарный делитель $f$, то его корни в $\CC$~--- некоторые из $\alpha_i$, а его
коэффициенты (с точностью до знака)~--- элементарные симметрические многочлены от его корней. Поэтому модули
коэффициентов $g$ ограничены через степень и коэффициенты $f$ (\rf{1.17b}). Кроме того, они целые по \rf{1.14}. Значит,
кандидатов в делители лишь конечное число, и для каждого делимость проверяется делением с остатком в $\Z[X]$ (\rf{1.14a}).
Если $f$ приводим, у него есть унитарный делитель степени $l$ с $1\le l\le m/2$; найдя его, применяем ту же процедуру
к $g$ и к $f/g$. Так за конечное число шагов получается разложение на неприводимые множители.
\end{remark}

Оценка корней на шаге 2 получается так: при большом $|x|$ старший член $x^m$ <<перевешивает>> сумму остальных, и многочлен
не может обратиться в нуль; нужно лишь понять, начиная с какого $|x|$ это гарантированно происходит.

\begin{proposition}[1.17a: оценка корней]
Пусть $f=X^m+a_1X^{m-1}+\dots+a_m\in\CC[X]$, $m\ge1$, и
\[
B=\max_{1\le i\le m}|a_i|,\qquad L=|a_1|+\dots+|a_m| .
\]
Тогда каждый корень $\alpha$ многочлена $f$ удовлетворяет неравенствам
\[
|\alpha|\le\max\{1,L\}\le\max\{1,mB\}\qquad\text{и}\qquad|\alpha|<1+B .
\]
\end{proposition}

\begin{proof}
Если $|\alpha|\le1$, то первое неравенство очевидно, а второе верно, так как либо $B>0$, либо $f=X^m$ и $\alpha=0$.
Пусть $|\alpha|>1$. Из $f(\alpha)=0$ имеем $\alpha^m=-(a_1\alpha^{m-1}+\dots+a_m)$, и по неравенству треугольника
\[
|\alpha|^m\le|a_1||\alpha|^{m-1}+|a_2||\alpha|^{m-2}+\dots+|a_m| .
\]
Так как $|\alpha|>1$, каждая степень $|\alpha|^{m-k}$, $1\le k\le m$, не больше $|\alpha|^{m-1}$. Поэтому
$|\alpha|^m\le L|\alpha|^{m-1}$, т.\,е.\ $|\alpha|\le L\le mB$. (Именно здесь нужно $|\alpha|>1$: при $|\alpha|<1$
младшие степени больше старших, отсюда и единица в максимуме.)

Пусть по-прежнему $|\alpha|>1$; тогда $B>0$ (иначе $f=X^m$ и $\alpha=0$). Если не огрублять сумму, а оценить её
геометрической прогрессией, то
\[
|\alpha|^m\le B\bigl(|\alpha|^{m-1}+\dots+|\alpha|+1\bigr)=B\,\frac{|\alpha|^m-1}{|\alpha|-1}<\frac{B\,|\alpha|^m}{|\alpha|-1},
\]
откуда $|\alpha|-1<B$.
\end{proof}

При $m\ge2$ и $B\ge1$ (так, например, для любого целочисленного $f\ne X^m$ степени $\ge2$) вторая оценка (оценка Коши)
не хуже грубой оценки $\max\{1,mB\}$: $1+B\le2B\le mB$. (При $m=1$ это не так: для $X+5$ имеем $1+B=6>5=mB$.) Но и
оценка через $L$ бывает лучше: для $X^4+4$ имеем $L=4<5=1+B$. Поэтому в качестве границы модулей корней удобно брать
$\rho=\min\bigl(\max\{1,L\},\,1+B\bigr)$.

\begin{proposition}[1.17b: оценка коэффициентов делителя]
Пусть все корни унитарного многочлена $f\in\CC[X]$ по модулю не больше $\rho$ (например, $\rho=1+B$), и
$g=X^l+b_1X^{l-1}+\dots+b_l$~--- унитарный делитель $f$ в $\CC[X]$. Тогда
\[
|b_j|\le\binom lj\rho^j\qquad(1\le j\le l).
\]
\end{proposition}

\begin{proof}
Имеем $g=\prod_{i=1}^l(X-\beta_i)$, где $\beta_i$~--- корни $f$ (с учётом кратности), и $b_j=(-1)^je_j(\beta_1,\dots,\beta_l)$,
где $e_j$~--- сумма всех $\binom lj$ произведений $j$ различных по номеру $\beta_i$. Каждое такое произведение по модулю не
больше $\rho^j$.
\end{proof}

Если $f\in\Z[X]$, то по \rf{1.14} коэффициенты $b_j$ целые, т.\,е.\ $b_j$~--- целое число из отрезка
$\bigl[-\binom lj\rho^j,\binom lj\rho^j\bigr]$. Поэтому кандидатов в унитарные делители степени $l$ не больше
\[
\prod_{j=1}^l\Bigl(2\Bigl\lfloor\binom lj\rho^j\Bigr\rfloor+1\Bigr)\le(2G+1)^l,\qquad G=\max_{1\le j\le l}\binom lj\rho^j .
\]

\begin{example}[1.17c]
Для $f=X^4+4$ имеем $B=L=4$, и можно взять $\rho=4$. Унитарный квадратный делитель $X^2+b_1X+b_2$ обязан иметь
$|b_1|\le8$, $|b_2|\le16$: это $17\cdot33=561$ кандидат (по оценке Коши $\rho=5$ получилось бы $21\cdot51=1071$
кандидатов); линейных делителей нет: $\pm1,\pm2,\pm4$ не корни. Перебор находит
$X^4+4=(X^2-2X+2)(X^2+2X+2)$.
\end{example}

Практически так, конечно, не делают: разложение в $\Q[X]$ и в $\F_p[X]$ умеют выполнять системы компьютерной алгебры
(например, PARI). Они дают, например,
\[
6X^2+18X-24=6(X-1)(X+4)\ \text{в }\Q[X],\qquad X^2+3X+3=(X+4)(X+6)\ \text{в }\F_7[X]
\]
(проверка: $(X+4)(X+6)=X^2+10X+24\equiv X^2+3X+3\pmod7$). Об устройстве быстрых алгоритмов~--- в отступлении \rf{1.18f}.

\subsection*{Приведение по модулю простого числа}

\begin{remark}[1.18]
Пусть $f\in\Z[X]$, а простое $p$ не делит старший коэффициент $f$. Нетривиальное разложение $f=gh$ в $\Z[X]$ (на
множители меньшей степени) даёт нетривиальное разложение $\bar f=\bar g\bar h$ в $\F_p[X]$: так как $p$ не делит
произведение старших коэффициентов $g$ и $h$, он не делит ни один из них, и $\deg\bar g=\deg g$, $\deg\bar h=\deg h$.
Вместе с леммой Гаусса (\rf{1.13}) это даёт: \emph{если $\bar f$ неприводим в $\F_p[X]$ для некоторого простого $p$, не
делящего старший коэффициент $f$, то $f$ неприводим в $\Q[X]$} (и в $\Z[X]$, если $f$ примитивен).

Условие на старший коэффициент нужно: $2X^2+X=X(2X+1)$, хотя $\overline{2X^2+X}=X$ неприводим в $\F_2[X]$.
\end{remark}

\begin{example}[1.18a]
(a) Многочлен $X^3-3X-1$ из \rf{1.12} по модулю $2$ равен $X^3+X+1$, у которого нет корней в $\F_2$; кубический
многочлен без корней неприводим, и мы снова видим, что $X^3-3X-1$ неприводим над $\Q$.

(b) $f=X^4+2X^3+3X+5$ неприводим над $\Q$: $\bar f=X^4+X+1\in\F_2[X]$ не имеет корней ($\bar f(0)=\bar f(1)=1$), а
единственный неприводимый квадратный многочлен над $\F_2$~--- это $X^2+X+1$, и $(X^2+X+1)^2=X^4+X^2+1\ne\bar f$.

(c) Иногда помогает сравнение нескольких простых. Пусть $f=X^4-3X^2-3X+1$. Тогда
\[
f\equiv(X+1)(X^3+X^2+1)\pmod2,\qquad f\equiv X^4+1\equiv(X^2+X+2)(X^2+2X+2)\pmod3,
\]
и все множители справа неприводимы (линейный~--- очевидно, а многочлены степени $2$ и $3$ не имеют корней в $\F_2$,
соответственно в $\F_3$). Если бы $f$ был приводим над
$\Q$, то по \rf{1.14} $f=gh$ с унитарными $g,h\in\Z[X]$, и (по однозначности разложения в $\F_p[X]$) $\bar g$ был бы
произведением некоторых из неприводимых множителей $\bar f$. Делитель степени $1$ или $3$ невозможен по модулю $3$, а
делитель степени $2$~--- по модулю $2$. Значит, $f$ неприводим, хотя по модулю $2$ и по модулю $3$ он приводим.
\end{example}

Этот признак очень полезен, но не всегда работает: многочлен может быть неприводим над $\Q$ и приводим по модулю
\emph{каждого} простого числа. Нам понадобится простой факт о квадратах в $\F_p$.

\begin{lemma}[1.18b]
Пусть $p$~--- нечётное простое. Квадраты образуют подгруппу индекса $2$ в $\F_p^\times$; следовательно, произведение двух
неквадратов из $\F_p^\times$~--- квадрат. В $\F_2$ каждый элемент~--- квадрат.
\end{lemma}

\begin{proof}
Отображение $\varphi\colon\F_p^\times\to\F_p^\times$, $x\mapsto x^2$,~--- гомоморфизм, его образ~--- подгруппа квадратов,
а ядро $\{x\mid x^2=1\}=\{1,-1\}$ (у $X^2-1$ не более двух корней по \rf{1.7}, и $1\ne-1$, так как $p\ne2$). Поэтому
подгруппа квадратов имеет порядок $(p-1)/2$ и индекс $2$, и факторгруппа $\F_p^\times/\operatorname{Im}\varphi$ имеет
порядок $2$. Неквадраты~--- это неединичный класс, а произведение двух неединичных элементов группы порядка $2$
единично. В $\F_2$: $0=0^2$, $1=1^2$.
\end{proof}

(То же следует из цикличности группы $\F_p^\times$, упражнение 1-3: квадраты~--- это чётные степени образующего.)

\begin{example}[1.18c]
Многочлен $f=X^4-10X^2+1$ неприводим в $\Z[X]$, но приводим по модулю каждого простого $p$.

\emph{Неприводимость над $\Q$.} По \rf{1.11} рациональные корни могут быть только $\pm1$, а $f(\pm1)=-8\ne0$; поэтому
линейных множителей нет. Если $f$ раскладывается в произведение двух квадратных, то по \rf{1.14} их можно взять
унитарными с целыми коэффициентами:
\[
X^4-10X^2+1=(X^2+aX+b)(X^2+cX+d),\qquad a,b,c,d\in\Z .
\]
Сравнивая коэффициенты при $X^3$, $X$, $X^2$ и $X^0$, получаем $c=-a$, $a(d-b)=0$, $b+d-a^2=-10$, $bd=1$. Из $bd=1$:
$b=d=\pm1$, и тогда $a^2=10\pm2\in\{8,12\}$, что не квадрат целого числа. Итак, $f$ неприводим над $\Q$ (и в $\Z[X]$,
так как он унитарен).

\emph{Приводимость по модулю $p$.} Следующие тождества верны в любом коммутативном кольце, содержащем элемент с
указанным квадратом:
\begin{align*}
X^4-10X^2+1&=(X^2-1)^2-8X^2=(X^2-2\sqrt2X-1)(X^2+2\sqrt2X-1),\\
X^4-10X^2+1&=(X^2+1)^2-12X^2=(X^2-2\sqrt3X+1)(X^2+2\sqrt3X+1),\\
X^4-10X^2+1&=(X^2-5)^2-24=\bigl(X^2-(5+2\sqrt6)\bigr)\bigl(X^2-(5-2\sqrt6)\bigr).
\end{align*}
Если $2$~--- квадрат в $\F_p$, работает первое; если $3$~--- квадрат в $\F_p$, то второе. Если ни $2$, ни $3$ не квадраты
в $\F_p$, то $p$ нечётно, $2,3\in\F_p^\times$, и по \rf{1.18b} $6=2\cdot3$~--- квадрат, так что работает третье. Например,
\[
f\equiv(X+1)^4\pmod2,\qquad f\equiv(X^2+1)^2\pmod3,\qquad f\equiv(X^2+2)(X^2+3)\pmod5,
\]
\[
f\equiv(X^2+X+6)(X^2+6X+6)\pmod7,\qquad f\equiv(X+2)(X+11)(X+12)(X+21)\pmod{23}.
\]
\end{example}

Корни $f$~--- четыре числа $\pm\sqrt2\pm\sqrt3$, а $f$~--- минимальный многочлен числа $\sqrt2+\sqrt3$. Три выписанных
разложения (над $\CC$)~--- это в точности три разбиения четырёх корней на пары. Отсюда ещё одно доказательство того,
что у $f$ нет квадратных множителей в $\Q[X]$: унитарный квадратный делитель имел бы вид $(X-r)(X-r')$ для двух корней
$r\ne r'$ (однозначность разложения в $\CC[X]$), а частное~--- вид $(X-r'')(X-r''')$ для двух других, т.\,е.\ получилось бы
одно из трёх выписанных разложений; но в каждом из них коэффициенты $2\sqrt2$, $2\sqrt3$, $5\pm2\sqrt6$ иррациональны
(\rf{1.11}).

Объяснение явления даёт
теория Галуа: группа Галуа $f$ изоморфна $(\Z/2\Z)^2$, в ней нет элемента порядка $4$, а по теореме Дедекинда \rf{4.28}
для $p$, не делящего дискриминант $2^{14}3^2$, степени неприводимых множителей $\bar f$~--- длины циклов некоторого
элемента группы Галуа, действующего на корнях; поэтому для таких $p$ многочлен $\bar f$ не может быть неприводимым. Общее изучение таких
многочленов требует неэлементарных методов; известно, например, что для каждого составного $n$ существует многочлен
степени $n$ из $\Z[X]$, неприводимый над $\Z$, но приводимый по модулю всех простых (Р.~Брандл, 1986). Похожее явление
для корней разобрано в упражнении 1-8: многочлен может иметь корень по модулю любого целого числа, не имея целых корней.

\begin{remark}[1.18d]
В частности, многочлен степени $4$ может не иметь корней в данном поле и всё же быть приводимым: так ведут себя, скажем, $\bar f$ по модулю $5$ и $7$ в примере \rf{1.18c}, а над $\Q$~---
$X^4+4=(X^2-2X+2)(X^2+2X+2)$ из \rf{1.17c}. Поэтому для многочленов степени $\ge4$ отсутствие корней не означает
неприводимости.
\end{remark}

\subsection*{Методы разложения многочленов}

\begin{aside}[1.18e: подъём Гензеля]
Пусть $f\in\Z[X]$ унитарен, $p$ простое, и $f\equiv gh\pmod{p^j}$ ($j\ge1$), где $g,h\in\Z[X]$ унитарны, а $\bar g,\bar h$
взаимно просты в $\F_p[X]$. Тогда найдутся унитарные $g'\equiv g$, $h'\equiv h\pmod{p^j}$ тех же степеней с
$f\equiv g'h'\pmod{p^{j+1}}$. Действительно, запишем $f-gh=p^je$, $e\in\Z[X]$. По \rf{1.8} в $\F_p[X]$ есть $\sigma,\tau$
с $\sigma\bar g+\tau\bar h=1$; поэтому сравнение $ah+bg\equiv e\pmod p$ разрешимо, причём можно взять $\deg a<\deg g$
(заменим $\tau\bar e$ остатком от деления на $\bar g$, как в доказательстве \rf{1.8}; тогда $\deg b<\deg h$ при
$\deg\bar e<\deg f$, что выполнено: $f$ и $gh$ унитарны и имеют одну степень~--- если бы, скажем, $\deg f>\deg gh$, то
старший коэффициент $1$ многочлена $f-gh$ делился бы на $p^j$,~--- поэтому $\deg e<\deg f$). Поднимем $a,b$ до многочленов из
$\Z[X]$ тех же степеней (коэффициенты~--- любые представители вычетов) и положим $g'=g+p^ja$, $h'=h+p^jb$; так как
$\deg a<\deg g$, $\deg b<\deg h$, многочлены $g',h'$ унитарны тех же степеней:
\[
g'h'=gh+p^j(ah+bg)+p^{2j}ab\equiv gh+p^je=f\pmod{p^{j+1}} .
\]
Повторяя, поднимаем разложение по модулю $p$ до разложения по модулю любой степени $p^k$.

Подъём единствен: если $f\equiv g'h'\equiv g''h''\pmod{p^k}$, где все четыре многочлена унитарны, $g'\equiv g''$,
$h'\equiv h''\pmod p$ и $\bar g',\bar h'$ взаимно просты, то $g'\equiv g''$, $h'\equiv h''\pmod{p^k}$. Действительно,
пусть сравнения $g'\equiv g''$, $h'\equiv h''$ верны по модулю $p^j$, $j<k$; запишем $g''=g'+p^ja$, $h''=h'+p^jb$ с
$a,b\in\Z[X]$, $\deg a<\deg g'$, $\deg b<\deg h'$ (из $\bar g'=\bar g''$ следует, что унитарные $g',g''$ имеют одну
степень, и в разности старшие члены сокращаются; так же для $h',h''$). Тогда
\[
0\equiv g''h''-g'h'=p^j(ah'+bg')+p^{2j}ab\equiv p^j(ah'+bg')\pmod{p^{j+1}},
\]
т.\,е.\ $\bar a\bar h'=-\bar b\bar g'$ в $\F_p[X]$. Так как $\bar g'$ и $\bar h'$ взаимно просты, $\bar g'\mid\bar a$, а
$\deg\bar a<\deg\bar g'$; значит, $\bar a=0$, и тогда $\bar b\bar g'=0$, $\bar b=0$. Итак, сравнения верны по модулю
$p^{j+1}$.
\end{aside}

\begin{aside}[1.18f: обзор методов разложения в {$\Z[X]$}]
Перечислим без доказательств основные подходы к разложению многочленов из $\Z[X]$ (а значит, по \rf{1.17}, и из
$\Q[X]$).
\begin{enumerate}[leftmargin=*,label=\arabic*.]
\item \emph{Перебор по оценке коэффициентов} (\rf{1.17}). Конечен, но число кандидатов растёт экспоненциально по степени
и по числу цифр коэффициентов.

\item \emph{Метод Кронекера.} Если $g\mid f$ в $\Z[X]$, то $g(k)\mid f(k)$ для каждого $k\in\Z$. Унитарный $g$ степени $l$
однозначно определяется значениями в $l$ различных точках $k_1,\dots,k_l$ (многочлен $g-X^l$ степени $<l$ восстанавливается
интерполяцией Лагранжа). Выберем $k_i$ с $f(k_i)\ne0$; тогда для каждого $i$ значение $g(k_i)$~--- один из конечного числа
делителей $f(k_i)$. Перебирая наборы делителей, строим кандидата, проверяем, что его коэффициенты целые, и делим.
(Для неунитарных многочленов берут $l+1$ точку.) Например, для $f=X^4+4$ и квадратного $g$ берём $k=0,1$: $g(0)\mid4$,
$g(1)\mid5$, т.\,е.\ $6\cdot4=24$ кандидата $X^2+(v-u-1)X+u$ с $u=g(0)$, $v=g(1)$; ровно два из них делят $f$:
$X^2+2X+2$ ($u=2$, $v=5$) и $X^2-2X+2$ ($u=2$, $v=1$).

\item \emph{Приведение по модулю $p$} (\rf{1.18}, \rf{1.18a}). Быстро доказывает неприводимость во многих случаях, но
не во всех (\rf{1.18c}).

\item \emph{Алгоритм Берлекэмпа--Цассенхауса} для унитарного $f\in\Z[X]$:
\begin{itemize}[leftmargin=*]
\item Избавиться от кратных множителей: в характеристике $0$ многочлен $f/\gcd(f,f')$ имеет те же неприводимые множители,
что и $f$, но без кратностей (НОД находится алгоритмом Евклида, \rf{1.8}). Далее $f$ без кратных множителей.
\item Выбрать простое $p$, по модулю которого $\bar f$ тоже без кратных множителей (годятся все $p$, кроме конечного
числа: по \rf{1.8f} достаточно, чтобы $p$ не делил целое $R\ne0$ с $Af+Bf'=R$, $A,B\in\Z[X]$; с точностью до знака $R$~---
дискриминант $f$).
\item Разложить $\bar f$ в $\F_p[X]$ (алгоритм Берлекэмпа). Если $\bar f=f_1\cdots f_r$ с различными неприводимыми $f_i$, то
по китайской теореме об остатках $A=\F_p[X]/(\bar f)\cong\prod_i\F_p[X]/(f_i)$, и каждый сомножитель~--- поле. В поле
характеристики $p$ уравнение $a^p=a$ имеет ровно $p$ решений~--- элементы простого подполя (их $p$, все они корни $X^p-X$
по малой теореме Ферма, а больше $p$ корней быть не может, \rf{1.7}). Поэтому $\F_p$-линейное подпространство
$\{a\in A\mid a^p=a\}$ имеет размерность $r$ и находится решением линейной системы. Для любого его элемента $v$,
представленного многочленом степени $<\deg\bar f$, имеем $v^p-v=\prod_{s\in\F_p}(v-s)\equiv0\pmod{\bar f}$, и множители
$v-s$ попарно взаимно просты; поэтому $\bar f=\prod_{s\in\F_p}\gcd(\bar f,v-s)$ (с точностью до константы), и это разложение
нетривиально, если $v$ непостоянен. Повторяя, получаем все $f_i$.
\item Поднять разложение до модуля $p^k$ (\rf{1.18e}, применённое последовательно к $f_1\cdot(f_2\cdots f_r)$,
затем к $f_2\cdot(f_3\cdots f_r)$ и т.\,д.): $f\equiv g_1\cdots g_r\pmod{p^k}$ с унитарными $g_i\in\Z[X]$, $\bar g_i=f_i$, где $p^k$ больше удвоенной оценки коэффициентов делителей
(\rf{1.17b}; на практике~--- более точной оценки Миньотта: если $X^l+b_1X^{l-1}+\dots+b_l$ делит $f$ в $\Z[X]$, то
$|b_j|\le\binom lj\|f\|_2$, где $\|f\|_2$~--- корень из суммы квадратов коэффициентов $f$). Тогда каждый делитель $g$ многочлена $f$ в $\Z[X]$
восстанавливается по своему вычету по модулю $p^k$, если брать вычеты из интервала $(-p^k/2,p^k/2)$.
\item Рекомбинация: каждый неприводимый множитель $f$ над $\Z$ по модулю $p^k$ равен произведению некоторого подмножества
поднятых множителей (если $G$~--- такой множитель, то $\bar G=\prod_{i\in S}f_i$ для некоторого $S$, и по
единственности подъёма \rf{1.18e}, применённой к разложениям $f=G\cdot(f/G)$ и
$f\equiv\prod_{i\in S}g_i\cdot\prod_{i\notin S}g_i$, получаем $G\equiv\prod_{i\in S}g_i\pmod{p^k}$). Перебираем подмножества по возрастанию размера и проверяем делимость в $\Z[X]$.
\end{itemize}
Все шаги, кроме последнего, выполняются за полиномиальное время; рекомбинация в худшем случае требует перебора $2^r$
подмножеств. Худший случай дают \emph{многочлены Суиннертона-Дайера}~--- минимальные многочлены сумм
$\sqrt{p_1}+\dots+\sqrt{p_n}$ для различных простых $p_i$ (при $n=2$, $p_1=2$, $p_2=3$ это $X^4-10X^2+1$). Их степень
равна $2^n$, а по модулю каждого простого они распадаются на множители степени не выше $2$, т.\,е.\ не менее чем на
$2^{n-1}$ множителей.

\item \emph{Полиномиальные алгоритмы.} Первый алгоритм разложения в $\Q[X]$ с полиномиальным временем работы построили
А.~Ленстра, Х.~Ленстра и Л.~Ловас (1982) с помощью редукции решёток (LLL). Современные реализации комбинируют
Берлекэмпа--Цассенхауса с редукцией решёток на шаге рекомбинации (М.~ван Хуй (M.~van~Hoeij), 2002).
\end{enumerate}
\end{aside}

Подъём Гензеля \rf{1.18e} поднимает разложения; для корней есть более простой вариант, годный для любого многочлена из
$\Z[X]$.

\begin{lemma}[1.18g: лемма Гензеля для корней]
Пусть $g\in\Z[X]$, $p$ простое, $x_1\in\Z$, $g(x_1)\equiv0\pmod p$ и $p\nmid g'(x_1)$. Тогда для каждого $k\ge1$ найдётся
$x_k\in\Z$ с
\[
x_k\equiv x_1\pmod p,\qquad g(x_k)\equiv0\pmod{p^k},
\]
причём такой $x_k$ единствен по модулю $p^k$ (в частности, $x_{k+1}\equiv x_k\pmod{p^k}$).
\end{lemma}

\begin{proof}
\emph{Существование}~--- индукцией по $k$. Пусть $x_k$ найден и $g(x_k)=p^ku$; ищем $x_{k+1}=x_k+p^ks$. По формуле Тейлора
(см.\ после \rf{1.7a})
\[
g(x_k+p^ks)\equiv g(x_k)+p^ks\,g'(x_k)=p^k\bigl(u+s\,g'(x_k)\bigr)\pmod{p^{2k}},
\]
а $2k\ge k+1$. Так как $x_k\equiv x_1\pmod p$, то $g'(x_k)\equiv g'(x_1)\not\equiv0\pmod p$, и сравнение
$u+s\,g'(x_k)\equiv0\pmod p$ разрешимо ($\F_p$~--- поле); при таком $s$ имеем $g(x_{k+1})\equiv0\pmod{p^{k+1}}$.

\emph{Единственность.} Пусть $y\equiv z\equiv x_1\pmod p$ и $g(y)\equiv g(z)\equiv0\pmod{p^k}$. Если уже $z=y+p^js$ с
$1\le j<k$, то, как выше, $0\equiv g(z)-g(y)\equiv p^js\,g'(y)\pmod{p^{j+1}}$, и так как $p\nmid g'(y)$, то $p\mid s$, т.\,е.\
$z\equiv y\pmod{p^{j+1}}$. По индукции $z\equiv y\pmod{p^k}$.
\end{proof}

Условие $p\nmid g'(x_1)$ означает, что $\bar x_1$~--- простой корень $\bar g$: $\bar g=(X-\bar x_1)\bar h$ с
$\bar h(\bar x_1)=\bar g'(\bar x_1)\ne0$, т.\,е.\ с взаимно простыми множителями, как в \rf{1.18e}. Без него лемма неверна:
у $X^2-13$ есть корень $0$ по модулю $13$, но сравнение $x^2\equiv13\pmod{169}$ неразрешимо (из $13\mid x$ следует
$169\mid x^2$). Для $g=X^2-a$ имеем $g'(x)=2x$, и условие выполнено, если $p$ нечётно и $p\nmid a$: тогда каждый квадратный
корень из $a$ по модулю $p$ поднимается до квадратного корня из $a$ по модулю любой степени $p$ (это используется в
упражнении 1-8).
